\section{The affine game supplied by the 2-to-1 Games Theorem}
\label{sec:affine-games}

The Max-Cut reduction needs two features beyond the usual numerical
2-to-1 game gap: the projections must be affine over $\F$, and the
alphabet must be fixed before the completeness error is chosen.
We derive these features from the construction and soundness analysis in
\cite[Sections~4--5]{DKKMS25}.  Its Grassmann agreement hypothesis is
supplied by the expansion theorem of Khot, Minzer, and Safra~\cite{KMS23}
and the expansion-to-agreement implications of Barak, Kothari, and
Steurer~\cite[Theorems~9--11]{BKS19}; the introduction of the published
version of~\cite{DKKMS25} records this completion.

We first describe the folded game and its affine projections.  The
source's vertex definition excludes certain subspaces that its displayed
sampler can produce.  We specify the resulting conditioning, complete the
tables used in its soundness proof, and give the invariant-list argument
with the needed exceptional-subspace estimates.  For the second player
we use bounded lists on quotient spaces, keeping each restriction to
the advice space separate.  Completeness then follows by choosing the
initial 3Lin error last.

\subsection{The folded game and its edge law}

\paragraph{Starting instance.}
Theorem~4.1 of~\cite{DKKMS25}, obtained from H\aa stad's
theorem~\cite{Hastad01}, supplies an absolute constant $s_*<1$ such that
for every sufficiently small $\varepsilon>0$ it is NP-hard to distinguish regular binary
3Lin instances of value at least $1-\varepsilon$ from those of value at
most $s_*$.  Regularity includes bounded variable occurrence, three
distinct variables per equation, and at most one variable shared by two
distinct equations.

For now let $\ell\ge2$, $k_0\ge\ell$, and $\beta\in(0,1)$ be fixed
parameters, with $\beta$ rational.  The two alphabet dimensions will be
$\ell$ and $\ell-1$, so the dimension in \Cref{prop:affine-games} is
$n=\ell-1$.  The soundness argument below chooses $\ell,k_0,\beta$
using only $b$ and the absolute constant $s_*$.

\paragraph{Spaces, labels, and folding.}
Let $\mathcal Z$ be the variables of the 3Lin instance, let
$V_0=\F^{\mathcal Z}$, and view each equation $e$ as
\[
  \langle x_e,\alpha\rangle=b_e,
  \qquad x_e\in V_0,\quad b_e\in\F,
\]
where $x_e$ has its three variable coordinates equal to one.  A
\emph{legitimate} tuple $U$ consists of $k_0$ disjoint equations with
the additional separation condition imposed in
\cite[Section~4.2]{DKKMS25}: variables from different equations of $U$
do not occur together in any input equation.  Put
\[
  V_U=\F^{\operatorname{var}(U)},
  \qquad H_U=\operatorname{span}\{x_e:e\in U\}.
\]
The vectors $x_e$, $e\in U$, are independent, so $\dim V_U=3k_0$,
$\dim H_U=k_0$, and there is a unique linear function
$h_U:H_U\to\F$ with $h_U(x_e)=b_e$.

Before folding, a left question is $(U,L)$, where
\[
  L\leq V_U,\qquad \dim L=\ell,\qquad L\cap H_U=\{0\}.
\]
For a variable subset $V\subseteq\mathcal Z$, write $V_V=\F^V\le V_0$;
the rule for sampling $V$ is given below. A right question is $(V,L')$,
with $L'\leq V_V$ and
$\dim L'=\ell-1$.  On an incident pair, $L'\leq L$.
The respective answers are linear functions on $L$ and $L'$, and the
constraint is restriction.  Choosing bases therefore identifies the
unfolded alphabets with $\F^\ell$ and $\F^{\ell-1}$.

The folding of~\cite[Section~4.2]{DKKMS25} puts $(U,L)$ and $(U_1,L_1)$
in the same class when
\[
 L+H_U+H_{U_1}=L_1+H_U+H_{U_1}.
\]
Lemma~4.2 there proves that this is an equivalence relation and supplies,
for every class $\mathcal C$,
a fixed $\ell$-dimensional space $R_{\mathcal C}$ such that
\begin{equation}
  R_{\mathcal C}+H_U=L+H_U
  \quad\text{for every }(U,L)\in\mathcal C.
  \label{eq:affine-complements}
\end{equation}
Both sides have dimension $\ell+k_0$; consequently
$R_{\mathcal C}\cap H_U=\{0\}$ as well.  A folded answer is a linear
function $\sigma:R_{\mathcal C}\to\F$.  Its unfolding on $(U,L)$ is
the restriction to $L$ of the unique extension
\begin{equation}
  \widetilde\sigma_U(r+h)=\sigma(r)+h_U(h),
  \qquad r\in R_{\mathcal C},\quad h\in H_U.
  \label{eq:affine-unfolding}
\end{equation}

For clarity, the affineness here is a property of the map between
\emph{answer spaces}.  For $z\in L$, write uniquely
$z=A_Uz+B_Uz$ with $A_Uz\in R_{\mathcal C}$ and $B_Uz\in H_U$.
The map $A_U:L\to R_{\mathcal C}$ is a linear isomorphism by
\eqref{eq:affine-complements}.  Thus unfolding is
\[
  \sigma\longmapsto\sigma\circ A_U+h_U\circ B_U,
\]
an affine bijection from $R_{\mathcal C}^*$ to $L^*$.
Composing it with the surjective restriction $L^*\to(L')^*$ gives a
surjective affine map with a one-dimensional kernel in its linear
part.  Every fiber has exactly two elements.  This proves the required
affine promise after bases are fixed.  Claim~4.4 of~\cite{DKKMS25}
shows that the projections agree when several unfolded edges are
aggregated into the same folded edge.  Equivalently, these edges may
be kept separately with their individual weights.

\paragraph{Validity of the edge sampling.}
We specify the treatment of ill-formed samples omitted from the displayed
sampler.  Sampling $k_0$ equations independently produces a
nonlegitimate tuple with probability
$O(k_0^2/|\mathcal Z|)$, by bounded occurrence.  For fixed $k_0$ this
can be made arbitrarily small.  If necessary, take sufficiently many
disjoint copies of the input instance; this preserves its value and
regularity.

Given a legitimate $U$, independently for each equation, the right
variable set $V$ contains all three variables with probability $1-\beta$
and otherwise one uniformly chosen variable.
If $f_0$ equations contribute all three, then
\[
  \dim V_V=k_0+2f_0,
  \qquad \dim(V_V\cap H_U)=f_0.
\]
For a uniformly chosen $(\ell-1)$-space $L'\leq V_V$, counting
nonzero vectors gives
\begin{align*}
  \Pr[L'\cap H_U\ne\{0\}]
  &\leq
  \frac{(2^{f_0}-1)(2^{\ell-1}-1)}{2^{k_0+2f_0}-1}
  \leq 2^{\ell-k_0}.
\end{align*}
If $L'\cap H_U=\{0\}$, a uniformly chosen $\ell$-space $L\leq V_U$
containing $L'$ is obtained by choosing a nonzero vector of $V_U/L'$.
It intersects $H_U$ nontrivially precisely when this vector belongs
to $(H_U+L')/L'$.  Hence
\[
  \Pr[L\cap H_U\ne\{0\}\mid L']
  =\frac{2^{k_0}-1}{2^{3k_0-\ell+1}-1}
  \leq 2^{\ell-2k_0}.
\]
We condition the sampling on legitimacy and $L\cap H_U=\{0\}$.
The discarded probability is at most
\begin{equation}
  \kappa=O(k_0^2/|\mathcal Z|)
          +2^{\ell-k_0}+2^{\ell-2k_0}.
  \label{eq:affine-invalid-probability}
\end{equation}
The two geometric terms can be made arbitrarily small by increasing
$k_0$ after $\ell$ is fixed.  We take their sum at most $1/4$ and,
by using disjoint copies of the input if needed, take the probability of
nonlegitimacy at most $1/4$.  Thus the probability $q_{\rm all}$ of
retaining a sample from the original independent-equation sampler is
at least $1/2$.  We distinguish this total retained probability from
the normalization after legitimacy has already been imposed.

Let $P$ be that sampler after conditioning on legitimate $U$, and put
$E=\{L\cap H_U=\{0\}\}$ and
$\eta=2^{\ell-k_0}+2^{\ell-2k_0}$.  Our edge law is $R=P(\,\cdot\mid E)$,
so $\TV(P,R)=P(E^c)\leq\eta$.  On the source's valid edge set define
the subprobability weights
\[
 w(e)=\sum_{\substack{\omega\in E\\\omega\mapsto e}}P(\omega),
 \qquad q_{\rm edge}=\sum_e w(e)=P(E)\geq1-\eta.
\]
Thus $R(e)=w(e)/q_{\rm edge}$. For any folded labeling, its acceptance
under $R$ is its raw valid acceptance divided by $q_{\rm edge}$. The raw
valid acceptance uses the law $P$ above, conditioned on legitimate $U$
but not on $E$, with payoff zero on invalid outcomes. We next bound this
quantity using completed tables in the published soundness analysis.

\subsection{Completed tables and the equations encoded by their lists}

The soundness proof uses tables on every $\ell$-dimensional subspace,
although a folded labeling specifies answers only at valid vertices.
We first extend those answers, preserving the consistency between the
tables on $V_U$ and $V_V$.  We then show that the linear functions
decoded from these tables still satisfy the equations in $U$.

\paragraph{Completing the tables.}
For a right variable set $V$ occurring with a legitimate $U$, define
\[
  H(V)=\operatorname{span}\{x_e:e\text{ is an input equation and }
                                     \operatorname{var}(e)\subseteq V\}.
\]
Every input equation supported in $\operatorname{var}(U)$ belongs to
$U$: an equation using two tuple blocks violates legitimacy, and a
distinct equation inside one block would share three variables with
that block's equation.  Since the equations of $U$ have disjoint
supports, this proves
\begin{equation}
  H_U\cap V_V=H(V).
  \label{eq:affine-validity-common-side-space}
\end{equation}
In particular, for $L\subseteq V_V$, validity of $(U,L)$ is independent
of the containing legitimate tuple $U$.

Fix a folded labeling.  Let $F_U[L]$ be its unfolded answer at each
valid $(U,L)$, and the zero linear function whenever
$L\cap H_U\ne\{0\}$.  For each $V$, define $F_V[L]$ using any
containing $U$ when $L\cap H(V)=\{0\}$, and by zero otherwise.
For valid $L$, all vertices $(U,L)$ with that same $L$ lie in one
folding class; their unfolded answers agree by
\cite[Claim~4.4]{DKKMS25}.  Equation~\eqref{eq:affine-validity-common-side-space}
gives the same independence for invalid $L$.  Hence the completed tables
satisfy
\begin{equation}
  F_U[L]=F_V[L]\qquad(L\subseteq V_V,\ \dim L=\ell).
  \label{eq:affine-validity-completed-consistency}
\end{equation}
The auxiliary table $F_V$ is on $\ell$-spaces; it is distinct from
the game's right-answer table on $(\ell-1)$-spaces.

The completed $F_U$ has a second property.  Fix
$K\supseteq H_U$ of dimension $\ell+k_0$.  All valid $L$ with
$L+H_U=K$ belong to one folding class.  The common folded label and
its $(H_U,h_U)$-extension give one linear function $f_K:K\to\F$
with $f_K|_{H_U}=h_U$ and $F_U[L]=f_K|_L$ for every such $L$.
This condition imposes no restriction on invalid entries.

\paragraph{Agreement and list parameters.}
The Grassmann agreement test for a table $F$ first chooses a uniform
$(\ell-1)$-space, then two independent uniform $\ell$-spaces containing
it, and accepts when their table entries agree on the common
$(\ell-1)$-space.  Write $\operatorname{agreement}(F)$ for its
acceptance probability.

Fix a soundness parameter $\delta\in(0,1)$.  The agreement lemma
\cite[Lemma~3.7]{DKKMS25} supplies constants
$q,r,C$ and a positive function $\alpha(\ell)$ at agreement threshold
$\delta^2/8$, as in the proof of its Lemma~5.3.  Set the thresholds
\[
  \tau_r=\min\{C/2,1/2\},\qquad
  \tau_i=10^{-9}\tau_{i+1}^{12}\quad(0\le i<r).
\]
Thus $0<\tau_i\le1/2$, and the pairs supplied by the agreement lemma
have agreement strictly above the relevant thresholds.
We also fix the second list thresholds
\[
 \zeta_r=\tau_0/2,\qquad
 \zeta_i=10^{-9}\zeta_{i+1}^{12}\quad(0\le i<r).
\]
Both vectors are fixed before choosing $\ell$;
the function $\alpha(\ell)$ is positive once $\ell$ is fixed.
We choose $\ell$ large enough for that agreement lemma and for the
ordinary list bound~\cite[Lemma~3.12]{DKKMS25} on
$(\ell-q)$-spaces with thresholds $(\zeta_i)$.  We also require
$q\le\ell-1$ and
\begin{equation}
  2^{q-\ell}\le\tau_0/4.
  \label{eq:affine-validity-ell-margin}
\end{equation}
The repetition parameter $k_0$ is chosen subsequently, large enough that
the cited ambient-dimension requirements hold for every dimension
between $k_0$ and $3k_0$ and for their quotients by a $q$-space.

For a fixed $q$-space $Q$, a list records linear functions on spaces
of codimension at most $r$ whose restrictions agree with the table on
more than the designated fraction of $\ell$-spaces containing $Q$.
Only pairs with no occurring extension are kept.  The next lemma
shows that, for the completed tables above, these maximal pairs
respect the original equations whenever $Q$ avoids their span.

\begin{lemma}[The invariant-list assertion on admissible advice spaces]
\label{lem:affine-validity-invariance}
Let $X$ have dimension $3k_0$, let $H\le X$ have dimension $k_0$,
and let $h:H\to\F$ be linear.  Suppose a table $F$ assigns a linear
function to every $\ell$-space, and that for every
$(\ell+k_0)$-space $K\supseteq H$ there is a linear function
$f_K:K\to\F$ such that $f_K|_H=h$ and
$F[L]=f_K|_L$ whenever $L+H=K$.
Fix a $q$-space $Q$ with $Q\cap H=\{0\}$, and set
\[
  e_{k_0}=2^{\ell+r-2k_0+1}.
\]
Assume $q\le\ell$, $3k_0-r>\ell$, \eqref{eq:affine-validity-ell-margin},
and $e_{k_0}\le\tau_0/4$.
For $Q\subseteq W\subseteq X$ of codimension $i\le r$, say that
$(g,W)$ occurs if $g:W\to\F$ is linear and
\[
  \Pr_{Q\subseteq L\subseteq W,\ \dim L=\ell}
       [F[L]=g|_L]>\tau_i.
\]
Call it maximal if it has no occurring extension to a strictly larger
space.  Every maximal occurring pair satisfies $H\subseteq W$ and
$g|_H=h$.
\end{lemma}

\begin{proof}
Put $H'=H\cap W$.  In $W/Q$, the image $(H'+Q)/Q$ has dimension
$\dim H'$ because $Q\cap H=\{0\}$.  Counting nonzero vectors in a
uniform $(\ell-q)$-space gives
\begin{align*}
 d_W&:=\Pr_{Q\subseteq L\subseteq W}[L\cap H\ne\{0\}]\\
 &\le
 \frac{(2^{\dim H'}-1)(2^{\ell-q}-1)}{2^{\dim W-q}-1}
 \le 2^{\ell+\dim H'-\dim W+1}
 \le e_{k_0}.
\end{align*}
The same bound holds with $W$ replaced by any containing space of
codimension at most $r$.

Partition the valid spaces $L$ with $Q\subseteq L\subseteq W$ by
$Y=L+H'$.  Within a fixed class, the table is the restriction of one
linear function on $Y$ respecting $h|_{H'}$; this follows by restricting
$f_{L+H}$.  If $g|_{H'}\ne h|_{H'}$, the difference of these two
functions is nonzero on $H'$.  A complement to $H'$ in $Y$ containing
$Q$ is a graph of a linear map on an $(\ell-q)$-dimensional quotient.
Each of its $\ell-q$ independent lift choices must satisfy one
nontrivial linear equation for the two functions to agree.
Thus agreement within this class is at most $2^{q-\ell}$ (and is zero
if the functions already disagree on $Q$).  Total agreement is at most
\[
  d_W+(1-d_W)2^{q-\ell}
  \le e_{k_0}+2^{q-\ell}\le\tau_0/2<\tau_i,
\]
a contradiction.  Hence $g|_{H'}=h|_{H'}$.

There is consequently a unique extension $\widetilde g$ to
$\widetilde W=W+H$ respecting $h$.  Under either of the two uniform
distributions on valid $\ell$-spaces containing $Q$ in $W$ and in
$\widetilde W$, the image of $L$ in
\[
  W/(W\cap H)\simeq\widetilde W/H
\]
is uniform among the $\ell$-spaces containing the image of $Q$.
Indeed, every such image has the same number of complements containing
$Q$, namely $2^{(\ell-q)\dim(H\cap W)}$ in $W$ and
$2^{(\ell-q)\dim H}$ in $\widetilde W$.
Moreover, agreement depends only on this image: $f_{L+H}$ and
$\widetilde g$ agree on $H$, so their difference vanishes on one
complement precisely when it vanishes on every complement.
The conditional agreement on valid spaces is therefore identical in
the two ambient spaces.  Since agreement in $W$ is at least $\tau_i$,
agreement in $\widetilde W$ is at least
\[
  (1-d_{\widetilde W})\frac{\tau_i-d_W}{1-d_W}
  \ge(1-e_{k_0})(\tau_i-e_{k_0})
  \ge\tau_i-2e_{k_0}\ge\tau_i/2.
\]
If $H\not\subseteq W$, the codimension $j$ of $\widetilde W$ is
strictly smaller than $i$, and $\tau_j\le\tau_{i-1}<\tau_i/2$.
This gives an occurring extension, contradicting maximality.
\end{proof}

Applying this lemma with $X=V_U$, $H=H_U$, and $h=h_U$ shows that
every maximal pair decoded from $F_U$ respects all equations in $U$,
provided $Q\cap H_U=\{0\}$.  We next estimate the probability of
that condition under the distribution used by the source's soundness
proof.

\subsection{Soundness, completeness, and the parameter order}

\paragraph{The advice distribution.}
The outer game in~\cite[Section~5.2]{DKKMS25} first samples $(U,V)$,
then takes $Q$ uniformly from the $q$-spaces of $V_V$.  It sends
$(U,Q)$ and $(V,Q)$ to its two players and asks for assignments on
$\operatorname{var}(U)$ and $V$, respectively.  It accepts when those
assignments agree on $V$ and the first satisfies every equation of $U$.
If $f_0$ equations contributed all three variables, then
$\dim V_V=k_0+2f_0$ and $\dim H(V)=f_0$.  Consequently, under this
actual distribution, uniformly in $U,V$,
\begin{equation}
 \Pr[Q\cap H_U\ne\{0\}\mid U,V]
 \le\frac{(2^{f_0}-1)(2^q-1)}{2^{k_0+2f_0}-1}
 \le2^{q-k_0}.
 \label{eq:affine-validity-advice-loss}
\end{equation}
We further increase $k_0$ until
\begin{equation}
 e_{k_0}\le\tau_0/4,
 \qquad 2^{q-k_0}\le\delta\alpha(\ell)/8.
 \label{eq:affine-validity-k-margins}
\end{equation}

\paragraph{From game acceptance to nonempty lists.}
Suppose a folded labeling has raw valid acceptance greater than
$\delta$ under $P$, the sampler conditioned on legitimate $U$ but not
on validity of $L$.  Testing the completed $F_U[L]$ against the
original right answer on $(V,L')$ can only increase this acceptance:
valid entries are unchanged, and invalid outcomes previously had payoff
zero.  We may therefore
apply the agreement/list and covering arguments in
\cite[Lemma~5.3, Claim~5.6, and Equation~(5.1)]{DKKMS25} to the
completed tables. Before invoking invariance, these arguments use total
linear tables, their agreement, and
\eqref{eq:affine-validity-completed-consistency}; all are now defined
on every sampled subspace.

For fixed $U,Q$, let $\mathcal L_{U,Q}$ be the uniform law on
$\ell$-spaces of $V_U$ containing $Q$.  Let $\mathcal L'_{U,Q}$
first sample $V$ conditional on the outer question $(U,Q)$ and then
a uniform $\ell$-space of $V_V$ containing $Q$.
These are the conditional laws in~\cite[Lemma~4.11]{DKKMS25}.
Indeed, if $N_q(d)$ denotes the number of $q$-spaces in $\F^d$, then
for $Q\subseteq V_V$ a uniform $\ell$-space $L\le V_V$ satisfies
\[
 \Pr[Q\subseteq L\mid V]=\frac{N_q(\ell)}{N_q(\dim V_V)}.
\]
Thus conditioning the source's $V,L$ sampler on $Q\subseteq L$
weights each $V$ proportionally to
$\Pr[V\mid U]/N_q(\dim V_V)$, exactly as in the outer game.

Let $E_1$ be the event
$\operatorname{agreement}(F_U)>\delta^2/8$, let $E_2$ be
\[
 \TV(\mathcal L_{U,Q},\mathcal L'_{U,Q})
 \le 2^{\ell+5}\sqrt\beta\,k_0^{1/4},
\]
and let $E_3$ be nonemptiness of the $Q$-list of $F_U$ at thresholds
$(\tau_i)$.  Event $E_2$ is precisely the smoothness condition in
the source's Equation~(4.3).
For $E_3$, Lemma~3.7 of the source supplies an occurring pair for an
$\alpha(\ell)$ fraction of uniform $Q$ whenever $E_1$ holds; any such
pair has a maximal occurring extension.  This uses only existence of
the first list, not a bound on its size.
The cited arguments give these events joint probability at least
$\delta\alpha(\ell)/4$.
Intersecting with $E_4=\{Q\cap H_U=\{0\}\}$ leaves probability at
least $\delta\alpha(\ell)/8$, by
\eqref{eq:affine-validity-advice-loss}--\eqref{eq:affine-validity-k-margins}.
The event $E_4$ depends only on $U,Q$.  Consequently it does not alter
the subsequent conditional law of $V$ after $U,Q$ have been fixed.

\paragraph{The second player's quotient lists.}
For each question $(V,Q)$ and each linear form $\gamma:Q\to\F$,
fix a linear extension $p_\gamma:V_V\to\F$.
Let $\nu:V_V\to V_V/Q$ be the quotient map.  Define a total table
$T_\gamma$ on the $(\ell-q)$-spaces of $V_V/Q$ by
\[
 T_\gamma[L/Q]=
 \begin{cases}
  \overline{F_V[L]-p_\gamma|_L},&F_V[L]|_Q=\gamma,\\
  0,&F_V[L]|_Q\ne\gamma.
 \end{cases}
\]
Here $Q\subseteq L\subseteq V_V$, $\dim L=\ell$, and the overline
denotes the induced linear form on $L/Q$; it is defined in the first
case because the difference vanishes on $Q$.

For this total table, use the ordinary maximal list with thresholds
$(\zeta_i)$: a pair $(g,S)$ occurs if $S\le V_V/Q$ has codimension
$i\le r$ and $T_\gamma$ agrees with $g$ on more than a $\zeta_i$
fraction of the $(\ell-q)$-spaces in $S$.
By~\cite[Lemma~3.12]{DKKMS25}, each such list has size at most
$M_0=M_0(\ell-q,r,\zeta_r)$, independent of $k_0$.
The second player chooses a uniform $\gamma\in Q^*$, then a uniform
pair $(g,S)$ from this list.  It lifts $g$ to the linear form
$g\circ\nu+p_\gamma$ on $\nu^{-1}(S)$, extends it uniformly to
$V_V$, and outputs the corresponding assignment.  If its chosen list
is empty, it uses an arbitrary assignment.  The entire procedure uses
only $(V,Q)$ and the fixed tables.

The zero entries may create additional occurring extensions in an
auxiliary table.  This causes no difficulty: the proof below needs
an extension of a specified linear form, rather than an identification
of these auxiliary maximal pairs with the maximal pairs of $F_V$.

\paragraph{Success of the two strategies.}
The first player chooses a uniform maximal pair $(f,W)$ from the
$\tau$-list of $F_U$ and extends $f$ uniformly to $V_U$, using an
arbitrary assignment if its list is empty.  This finite list need not
have a size bound independent of $k_0$: the success estimate will hold
for every one of its elements.  The second player uses the quotient
lists just constructed.

Fix $U,Q$ in $E_1\cap E_2\cap E_3\cap E_4$, and any first-player
choice $(f,W)$, with $i=\operatorname{codim}_{V_U}W\le r$.
By \Cref{lem:affine-validity-invariance}, $H_U\subseteq W$ and
$f|_{H_U}=h_U$.  Every extension chosen by the first player therefore
satisfies all equations of $U$.  This is the only use of the
invariant-list assertion in the source's decoding argument.

The covering step of~\cite[Lemmas~4.11 and~5.5 and the proof of
Lemma~5.3]{DKKMS25}, using
\eqref{eq:affine-validity-completed-consistency}, gives probability at
least $0.3\tau_i\ge0.3\tau_0$ over $V$ that
\[
 \operatorname{codim}_{V_V}(W\cap V_V)=i,
 \qquad
 \Pr_{Q\subseteq L\subseteq W\cap V_V}
       [F_V[L]=f|_L]>\tau_i/2.
\]
Here $L$ has dimension $\ell$.  Set $S_0=W\cap V_V$ and
$\gamma=f|_Q$.  Every equality $F_V[L]=f|_L$ in this display is
retained in $T_\gamma$.  Since
$\zeta_i\le\tau_0/2\le\tau_i/2$, the quotient pair
\[
 \bigl(\overline{f|_{S_0}-p_\gamma|_{S_0}},\ S_0/Q\bigr)
\]
occurs for $T_\gamma$ and hence has a maximal occurring extension
in that auxiliary table.  Lifting any such extension by the same
$p_\gamma$ gives a function agreeing with $f$ on all of $S_0$.
This conclusion concerns function extension and is unaffected by
the extra agreements that zero completion may introduce.

The second player selects this tag and one such maximal extension
with probability at least $1/M$, where
$M:=2^qM_0$ is independent of $k_0$.
Its eventual answer then agrees with $f$ on $W\cap V_V$.
The first player's uniform extension agrees with this
answer on all of $V_V$ with probability $2^{-i}\ge2^{-r}$:
the quotient $V_V/(W\cap V_V)$ has dimension $i$.
Thus, conditionally on $E_1\cap E_2\cap E_3\cap E_4$, the strategy
wins with probability at least
\[
 c_*:=2^{-r}\frac{0.3\tau_0}{M}>0.
\]
The overall success probability is at least
$\delta\alpha(\ell)c_*/8$, a positive constant independent of $k_0$.
The upper bound in~\cite[Lemma~5.4]{DKKMS25} tends to zero as
$\beta k_0\to\infty$.  Choosing $k_0$ sufficiently large gives the
same contradiction as in the source and proves that raw valid
acceptance is at most $\delta$.

\paragraph{Fixing soundness before completeness.}
For the game in \Cref{prop:affine-games}, take $\delta=b/4$ and retain
the valid raw weights.  Their total mass $q_{\rm edge}$ is at least
$1/2$ by the preceding rejection estimates.  Normalization therefore
gives
\[
  \Val\le\delta/q_{\rm edge}\le b/2<b.
\]
Every restriction above concerns only $b$ and parameters fixed from it.
After choosing $\ell$, take $k_0=j^4$ and $\beta=j^{-3}$ with $j$
sufficiently large.  Then $\beta k_0\to\infty$,
$\beta\sqrt{k_0}\to0$, and $\sqrt\beta\,k_0^{1/4}\to0$.
These limits satisfy the outer-game soundness and covering requirements,
including $2^\ell\beta\le1/8$, as well as all the displayed margins.
In particular, the alphabet dimension $n=\ell-1$ has been fixed using
only $b$.

\paragraph{Completeness with the alphabet fixed.}
Let $\alpha:\mathcal Z\to\F$ satisfy all but an $\varepsilon$ fraction of
the input equations, and also write $\alpha:V_0\to\F$ for its linear
extension.  Give the folded left question $\mathcal C$ the answer
$\alpha|_{R_{\mathcal C}}$ and give $(V,L')$ the answer $\alpha|_{L'}$.
Whenever every equation of sampled $U$ is satisfied,
$h_U=\alpha|_{H_U}$, so~\eqref{eq:affine-unfolding} is exactly
$\alpha|_{R_{\mathcal C}+H_U}$.  The edge is therefore satisfied.
This argument concerns the sampled edge; it does not require every
member of its folding class to involve satisfied equations.

Before either conditioning step, a union bound gives probability at most
$k_0\varepsilon$ of sampling an unsatisfied equation.  Since the
total retained probability $q_{\rm all}$ is at least $1/2$, the
conditioned game's value is at least
$1-k_0\varepsilon/q_{\rm all}\ge1-2k_0\varepsilon$.
After $\ell,k_0$ have been fixed for
soundness, choose $\varepsilon\leq a/(2k_0)$ in the initial 3Lin
gap.  This gives the asserted completeness without altering $n$.

\paragraph{Effectivity and size.}
All spaces in an individual question have bounded dimension.  There
are polynomially many $k_0$-tuples of input equations, and each tuple
has only constantly many subspaces, bases, and local sampling
outcomes.  The folding equivalence relation is explicitly testable
by linear algebra.  To find $R_{\mathcal C}$ effectively, start from
any $(U,L)\in\mathcal C$ and enumerate the $\ell$-spaces inside
$L+H_U$, whose dimension is the fixed constant $\ell+k_0$; test each
candidate against~\eqref{eq:affine-complements} for every class
member.  Lemma~4.2 guarantees a successful candidate.  Gaussian
elimination then computes bases and the affine edge maps.
The finite sampling distributions, their conditioning, and the
aggregation of parallel edges have rational weights of polynomial
bit length.  These observations give the deterministic polynomial
reduction claimed in the proposition.
This completes the proof of \Cref{prop:affine-games}.
